\documentclass[10pt,a4paper]{article}
\usepackage[a4paper,margin=14mm,top=12mm,bottom=14mm]{geometry}
\usepackage{fontspec}
\setmainfont{Liberation Sans}
\usepackage{polyglossia}
\setdefaultlanguage{french}
\usepackage{amsmath}
\usepackage{xcolor}
\usepackage{graphicx}
\usepackage[most]{tcolorbox}
\usepackage{tabularx,booktabs,array,multirow}
\usepackage{enumitem}
\usepackage{fancyhdr}
\usepackage{lastpage}
\usepackage[hidelinks]{hyperref}
\definecolor{uimm}{HTML}{005677}
\definecolor{cyan}{HTML}{13A2DD}
\definecolor{fond}{HTML}{E6F2FA}
\definecolor{ok}{HTML}{2E7D32}
\definecolor{ko}{HTML}{B3261E}
\definecolor{violet}{HTML}{7B1FA2}
\setlength{\parindent}{0pt}
\setlength{\parskip}{3pt}
\setlist{nosep}
\setlist[itemize]{leftmargin=13pt,label=\textcolor{cyan}{\textbullet}}
\newcommand{\titre}[1]{\vspace{5pt}{\color{uimm}\large\bfseries #1}\par\vspace{-3pt}{\color{cyan}\rule{\linewidth}{1.2pt}}\par\vspace{0pt}}
\tcbset{boite/.style={colback=fond,colframe=uimm,boxrule=0.6pt,arc=2pt,left=6pt,right=6pt,top=4pt,bottom=4pt},
        defi/.style={colback=violet!6,colframe=violet,boxrule=0.8pt,arc=2pt,left=6pt,right=6pt,top=4pt,bottom=4pt,
                     title={\bfseries Défi},coltitle=white,colbacktitle=violet,fonttitle=\small}}
\newcommand{\bandeau}[3]{%
\begin{tcolorbox}[colback=uimm,colframe=uimm,arc=2pt,left=4pt,right=8pt,top=4pt,bottom=4pt]
\begin{minipage}[c]{18mm}\setlength{\fboxsep}{1.2pt}\fcolorbox{white}{white}{\includegraphics[height=16.5mm]{logo_uimm.jpg}}\end{minipage}\hspace{4mm}%
\begin{minipage}[c]{\dimexpr\linewidth-22mm\relax}\color{white}
{\small #1}\\[1pt]{\LARGE\bfseries #2}\\[2pt]{\large #3}\hfill{\small CQPM Technicien qualité}
\end{minipage}
\end{tcolorbox}}
\newcommand{\reponse}[1][1]{\par\vspace{2pt}{\color{gray!60}\foreach \i in {1,...,#1}{\rule{\linewidth}{0.3pt}\par\vspace{9pt}}}}
\usepackage{pgffor}
\pagestyle{fancy}\fancyhf{}
\renewcommand{\headrulewidth}{0pt}
\newcommand{\pied}[1]{\fancyfoot[L]{\footnotesize\color{uimm}#1}\fancyfoot[R]{\footnotesize\color{uimm}Page \thepage{} sur \pageref{LastPage}}}

\usepackage{tikz}
\usepackage{longtable}
\renewcommand{\bandeau}[3]{%
\begin{tcolorbox}[colback=uimm,colframe=uimm,arc=2pt,left=4pt,right=8pt,top=4pt,bottom=4pt]
\begin{minipage}[c]{18mm}\setlength{\fboxsep}{1.2pt}\fcolorbox{white}{white}{\includegraphics[height=16.5mm]{logo_uimm.jpg}}\end{minipage}\hspace{4mm}%
\begin{minipage}[c]{\dimexpr\linewidth-22mm\relax}\color{white}
{\small #1}\\[1pt]{\LARGE\bfseries #2}\\[2pt]{\large #3}
\end{minipage}
\end{tcolorbox}}
% en-tête d'une fiche : code, titre, question, bloc, durée, organisation, niveau
\newcommand{\entete}[7]{%
\begin{tcolorbox}[colback=uimm,colframe=uimm,arc=2pt,left=4pt,right=6pt,top=4pt,bottom=4pt]
\begin{minipage}[c]{18mm}\setlength{\fboxsep}{1.2pt}\fcolorbox{white}{white}{\includegraphics[height=16.5mm]{logo_uimm.jpg}}\end{minipage}\hspace{3mm}%
\begin{minipage}[c]{\dimexpr\linewidth-58mm\relax}\color{white}\raggedright
{\small Base SPC · #4}\\[1pt]{\Large\bfseries #1 · #2}\\[2pt]{\itshape #3}
\end{minipage}\hfill
\begin{minipage}[c]{33mm}\color{white}\small\raggedleft
Durée : \textbf{#5}\\ \textbf{#6}\\ Niveau #7 sur 3
\end{minipage}
\end{tcolorbox}}
\newcommand{\tagE}{{\setlength{\fboxsep}{1.5pt}\colorbox{cyan}{\color{white}\scriptsize\bfseries Ellistat}}}
\newcommand{\tagW}{{\setlength{\fboxsep}{1.5pt}\colorbox{uimm}{\color{white}\scriptsize\bfseries Page web}}}
\tcbset{pred/.style={colback=orange!6,colframe=orange!80!black,boxrule=0.6pt,arc=2pt,left=6pt,right=6pt,top=3pt,bottom=1pt,
                    title={\bfseries Je prédis, avant de commencer},coltitle=white,colbacktitle=orange!80!black,fonttitle=\small},
        retenir/.style={colback=fond,colframe=uimm,boxrule=0.6pt,arc=2pt,left=6pt,right=6pt,top=3pt,bottom=3pt,
                    title={\bfseries À retenir},coltitle=white,colbacktitle=uimm,fonttitle=\small},
        plus/.style={colback=violet!6,colframe=violet,boxrule=0.6pt,arc=2pt,left=6pt,right=6pt,top=3pt,bottom=3pt,
                    title={\bfseries Pour aller plus loin},coltitle=white,colbacktitle=violet,fonttitle=\small},
        act/.style={colback=uimm!10,colframe=uimm,boxrule=0.5pt,arc=2pt,left=5pt,right=5pt,top=2pt,bottom=2pt,before skip=7pt,after skip=3pt},
        verif/.style={colback=orange!6,colframe=orange!80!black,boxrule=0.6pt,arc=2pt,left=6pt,right=6pt,top=3pt,bottom=3pt,
                    title={\bfseries À vérifier sur poste avant la première séance},coltitle=white,colbacktitle=orange!80!black,fonttitle=\small}}
\newcommand{\fichier}[1]{\texttt{\small #1}}
\newcommand{\menu}[1]{\textsf{\textbf{#1}}}
\newcommand{\champ}[1]{\rule{#1}{0.3pt}}

\pied{Base SPC · fiche M16 · Deux grades, un micromètre}
\begin{document}
\thispagestyle{fancy}
\entete{M16}{Deux grades, un micromètre}{Mon micromètre voit-il la différence entre une cale de grade 0 et une cale de grade 1 ?}{Le moyen de mesure}{30 min}{en binôme}{2}
\begin{tcolorbox}[boite]\textbf{Le but.} Mesurer deux cales de même cote et de grades différents, et voir ce que le micromètre en dit.\par\textbf{Il vous faut.} Une cale céramique de 10 mm, grade 0 · Une cale acier de 10 mm, grade 1 · Un micromètre 0-25 mm · Ellistat\end{tcolorbox}
\begin{tcolorbox}[boite]\textbf{Un grade, c'est quoi ?} La norme ISO 3650 définit quatre grades de cales : \textbf{K}, \textbf{0}, \textbf{1} et \textbf{2}. Chaque grade fixe deux tolérances : l'\textbf{écart limite} de la longueur par rapport à la cote marquée, et la \textbf{variation de longueur} sur les faces de mesure de la cale. Pour les grades 0, 1 et 2, plus le chiffre est petit, plus ces tolérances sont serrées. Le grade K est la cale de référence qui sert à étalonner les autres. Usages courants : grade 0, ateliers d'outillage, services de contrôle et réglage de comparateurs ; grade 1, contrôle de première pièce et vérification sur machine à mesurer tridimensionnelle ; grade 2, référence d'atelier et calage de butées de machine. Le certificat de vos cales fait foi.\par{\footnotesize Sources : ISO 3650 (édition 1998), résumé en ligne, \href{https://www.iso.org/standard/1241.html}{\nolinkurl{iso.org/standard/1241.html}} ; Aims Industrial, Gauge block guide (distributeur : usages et valeurs), \href{https://aimsindustrial.com.au/blogs/product-guides/gauge-block-guide}{\nolinkurl{aimsindustrial.com.au/blogs/product-guides/gauge-block-guide}}.}\end{tcolorbox}
\begin{tcolorbox}[pred]La cale de grade 0 est plus précise que celle de grade 1. Votre micromètre au micron le verra-t-il ? Écrivez l'écart que vous attendez entre les deux moyennes, en µm.\reponse\end{tcolorbox}
\titre{À faire}
\begin{enumerate}[label=\textcolor{cyan}{\textbf{\arabic*.}},leftmargin=15pt,itemsep=1.5pt]
\item Lire dans le tableau ci-dessous les écarts limites des deux grades pour une cale de 10 mm. Calculer l'écart maximal possible entre les deux cales : les deux limites s'ajoutent.
\item Nettoyer les deux cales. Mesurer \textbf{15 fois} chacune au cliquet, en alternant : céramique, acier, céramique, acier. Retirer la cale à chaque mesure.
\item \tagE\ \menu{Gage R\&R}, onglet \menu{Justesse} : deux études, « Ceramique\_grade0 » et « Acier\_grade1 ». Nombre de mesures 15 ; Valeur vraie 10,000 ; LSI 9,925 et LSS 10,075 (tolérance de la butée, $T = 0{,}150$ mm, centrée sur la cale : $10{,}000 \pm 0{,}075$) ; k 3. Relever Cg et Cgk.\par{\small \textbf{k} : nombre d'écarts types retenus de chaque côté de la moyenne ; $k=3$ donne une dispersion de $6s$. \textbf{Cg} compare cette dispersion à 20 \% de la tolérance : $C_g = 0{,}2\,T/(6s)$, avec $T$ = LSS $-$ LSI. \textbf{Cgk} retire en plus le décalage entre la moyenne et la valeur vraie $x_m$ : $C_{gk} = (0{,}1\,T - |\bar x - x_m|)/(3s)$. Le livret Bosch demande 1,33 au moins pour les deux.}
\item \tagW\ Coller les deux séries dans la page « Comparer deux séries de mesures », avec 10,000 comme valeur de la cale, à l'adresse suivante :\par\href{https://sjaubert.github.io/000_OI/SPC/Pages/Comparer_deux_series.html}{\nolinkurl{https://sjaubert.github.io/000_OI/SPC/Pages/Comparer_deux_series.html}}
\item Avec 15 mesures par cale, le plus petit écart de moyennes que le test déclare significatif vaut environ $0{,}75 \times s$. Prendre pour $s$ la moyenne des deux écarts-types.
\end{enumerate}
\vspace{3pt}{\renewcommand{\arraystretch}{1.35}\small\begin{tabularx}{\linewidth}{|l|*{3}{>{\centering\arraybackslash}X|}}\hline
\textbf{Écart limite de longueur (µm)} & \textbf{Grade 0} & \textbf{Grade 1} & \textbf{Grade 2}\\\hline
Cale jusqu'à 10 mm & $\pm 0{,}12$ & $\pm 0{,}20$ & $\pm 0{,}45$\\\hline
Cale de plus de 10 mm à 25 mm & $\pm 0{,}14$ & $\pm 0{,}30$ & $\pm 0{,}60$\\\hline
\end{tabularx}}\par{\footnotesize ISO 3650, valeurs relevées sur des pages de revendeurs : le certificat de vos cales fait foi.}\par\vspace{5pt}
\begin{minipage}[t]{0.485\linewidth}\textbf{Cale céramique grade 0 (mm)}\par\vspace{2pt}
{\renewcommand{\arraystretch}{1.4}\begin{tabularx}{\linewidth}{|*{3}{>{\centering\arraybackslash}p{5mm}|>{\centering\arraybackslash}X|}}\hline
\scriptsize 1 &  & \scriptsize 2 &  & \scriptsize 3 & \\\hline
\scriptsize 4 &  & \scriptsize 5 &  & \scriptsize 6 & \\\hline
\scriptsize 7 &  & \scriptsize 8 &  & \scriptsize 9 & \\\hline
\scriptsize 10 &  & \scriptsize 11 &  & \scriptsize 12 & \\\hline
\scriptsize 13 &  & \scriptsize 14 &  & \scriptsize 15 & \\\hline
\end{tabularx}}\end{minipage}\hfill\begin{minipage}[t]{0.485\linewidth}\textbf{Cale acier grade 1 (mm)}\par\vspace{2pt}
{\renewcommand{\arraystretch}{1.4}\begin{tabularx}{\linewidth}{|*{3}{>{\centering\arraybackslash}p{5mm}|>{\centering\arraybackslash}X|}}\hline
\scriptsize 1 &  & \scriptsize 2 &  & \scriptsize 3 & \\\hline
\scriptsize 4 &  & \scriptsize 5 &  & \scriptsize 6 & \\\hline
\scriptsize 7 &  & \scriptsize 8 &  & \scriptsize 9 & \\\hline
\scriptsize 10 &  & \scriptsize 11 &  & \scriptsize 12 & \\\hline
\scriptsize 13 &  & \scriptsize 14 &  & \scriptsize 15 & \\\hline
\end{tabularx}}\end{minipage}\par
\newpage
\titre{J'observe, je décide}
\begin{enumerate}[label=\textcolor{cyan}{\textbf{Q\arabic*.}},leftmargin=22pt]
\item Écart maximal possible entre les deux cales : \champ{14mm} µm. Écart observé entre vos deux moyennes : \champ{14mm} µm.\par\vspace{5pt}
\item La page conclut-elle à une différence réelle entre les deux moyennes ? Cela prouve-t-il que les deux cales sont identiques ?\reponse[2]
\item Plus petit écart détectable ($0{,}75 \times s$) = \champ{14mm} µm. Est-il plus petit ou plus grand que l'écart maximal possible entre les deux cales ?\reponse[1]
\item À quoi sert alors une cale de grade 0 ? Avec quel type d'instrument son grade devient-il visible ?\reponse[2]
\end{enumerate}
\begin{tcolorbox}[retenir]Un étalon doit être nettement meilleur que l'instrument qu'il qualifie. Plus fin que nécessaire, son grade ne se voit pas.\end{tcolorbox}
\begin{tcolorbox}[plus]Refaites avec 50 mesures par cale : le plus petit écart détectable tombe à environ $0{,}4 \times s$. Frôle-t-il l'écart maximal possible ? Que faudrait-il contrôler avant de conclure ?\end{tcolorbox}
\vfill{\footnotesize\color{uimm}\textbf{Ressources de la base :} Comparer deux séries de mesures · Données de secours.}
\end{document}